problem:  
Let $\mathbb{B}^3 \subset \mathbb{R}^3$ denote the unit ball with boundary $\mathbb{S}^2$.  
Consider a smooth foliation of $\mathbb{S}^2 \setminus \{p_*\}$ by embedded circles $\{\Gamma_t\}_{t\in[0,1)}$ that converge to a single point $p_* \in \mathbb{S}^2$ as $t\to1$.  
For each $t$, let $\mathcal{M}(\Gamma_t)$ denote the (possibly non-unique) collection of embedded area-minimizing disks in $\mathbb{B}^3$ with boundary $\Gamma_t$.

**Question:**  
Does there exist such a boundary foliation $\{\Gamma_t\}$ for which there is a corresponding family $\{M_t \in \mathcal{M}(\Gamma_t)\}_{t\in[0,1)}$ satisfying:

1. each $M_t$ is an embedded minimal disk in $\mathbb{B}^3$;  
2. the collection $\{M_t\}$ is disjoint and locally smooth except possibly at a limiting lamination singularity at the center of the ball; and  
3. $\{M_t\}$ forms a **minimal lamination** of $\mathbb{B}^3 \setminus \{0\}$, meaning that  
\[
\mathbb{B}^3 \setminus \{0\} = \bigcup_{t \in [0,1)} M_t,
\]  
and that $M_t$ converges (in the varifold or smooth lamination sense) to a limiting minimal cone or point singularity as $t \to 1$.

In other words: can $\mathbb{B}^3$ be laminated by embedded minimal disks whose boundaries form a smooth circle foliation on $\mathbb{S}^2$, degenerating only at the center?

Why is it a "good" problem:  
This formulation transforms the earlier heuristic foliation idea into a rigorously grounded problem in the modern framework of **minimal laminations and Plateau theory**. It avoids assuming smooth dependence of Plateau solutions, and instead asks whether such a global smooth or laminar family exists, thereby squarely addressing one of the central unresolved questions of geometric analysis: how solutions to the nonlinear minimal surface equation organize globally when boundary data vary.  

The problem lies at the intersection of **geometric measure theory**, **analysis of PDE**, and **low-dimensional topology**. Existence of a minimal lamination filling the ball would deeply connect boundary geometry with global interior regularity, and any counterexample would clarify limits of continuity for the Plateau problem.  

It is simple and visual (“Can all minimal disks spanning concentric-like circles on the sphere fill the ball as a lamination?”), yet it demands profound new understanding of compactness, uniqueness, and stability of minimal surfaces under varying boundaries. A resolution would likely require developing techniques bridging **boundary variation theory**, **lamination compactness**, and **singular foliation theory**, thus constituting a genuinely deep and fertile research problem.